% Graber--Pandharipande 1999 (Invent. Math. 135), "Localization of virtual classes"
% Edidin--Graham 1998, "Equivariant intersection theory"
% Kiem--Li 2013 (JAMS 26), "Localizing virtual cycles by cosections"
% Pandharipande--Thomas 2014 (Forum Math. Pi 2), "The Katz-Klemm-Vafa conjecture for K3 surfaces"
% Oberdieck--Pandharipande 2016 (arXiv:1411.1514), "Curve counting on K3 x E, the Igusa cusp form chi_10, and descendent integration"
% Maulik--Oberdieck--Pandharipande 2017 (arXiv:1605.05473), "Curves on K3 surfaces"
% Oberdieck--Pixton 2017 (arXiv:1706.10100), "Holomorphic anomaly equations for K3 surfaces"
% Oberdieck 2018, "Gromov-Witten invariants of the Hilbert schemes of points of a K3 surface"
% Nakajima 1997, "Heisenberg algebra and Hilbert schemes of points on projective surfaces" (Ann. Math. 145)
% Grojnowski 1996 (arXiv:alg-geom/9506020), "Instantons and affine algebras I"
% Behrend--Fantechi 1997, "The intrinsic normal cone"

\providecommand{\Coh}{\mathrm{Coh}}
\providecommand{\Hilb}{\mathrm{Hilb}}
\providecommand{\Db}{D^{b}}
\providecommand{\PT}{\mathrm{PT}}
\providecommand{\DT}{\mathrm{DT}}
\providecommand{\CoHA}{\mathrm{CoHA}}
\providecommand{\cO}{\mathcal O}
\providecommand{\cH}{\mathcal H}
\providecommand{\cF}{\mathcal F}
\providecommand{\bC}{\mathbb C}
\providecommand{\bQ}{\mathbb Q}
\providecommand{\bZ}{\mathbb Z}
\providecommand{\bP}{\mathbb P}
\providecommand{\bH}{\mathbb H}
\providecommand{\Gm}{\mathbb G_m}
\providecommand{\vir}{\mathrm{vir}}
\providecommand{\kch}{\kappa_{\mathrm{ch}}}
\providecommand{\kcat}{\kappa_{\mathrm{cat}}}
\providecommand{\kbkm}{\kappa_{\mathrm{BKM}}}
\providecommand{\kfib}{\kappa_{\mathrm{fiber}}}
\providecommand{\Deltafive}{\Delta_5}
\providecommand{\Dfive}{D_5}
\providecommand{\PhiOP}{\Phi_{10}^{\mathrm{OP}}}

\section*{Reduced equivariance on $K3\times E$}

Let \(S\) be a smooth projective K3 surface, let \(E\) be an elliptic
curve, and set \(X=S\times E\).  Since \(H^0(S,T_S)=0\),
\(\mathrm{Aut}(S)^\circ\) is trivial.  The connected automorphism group
of \(X\) is the translation group \(E\), and an abelian variety contains
no nontrivial affine algebraic subgroup \(\Gm\).  Ordinary
Atiyah--Bott or Graber--Pandharipande localisation would require an
additional global algebraic \(\Gm\)-action on the moduli problem.

The usable structures are narrower.  A local noncompact chart may carry
an \(\Omega\)-background \(\Gm\)-action.  A special K3 surface may carry
a finite symplectic automorphism.  The product always carries the
K3-sourced cosection of the stable-pair obstruction sheaf and the free
translation action of \(E\).  The compact enumerative theorem for
\(K3\times E\) is the reduced theory obtained from this cosection and the
quotient by the free elliptic translation.

The reduced DT statements below are in K3-primitive curve class unless a
nonprimitive divisor-sum correction is named.  The obstruction theory is
the Kiem--Li/Maulik--Pandharipande--Thomas reduced theory obtained from
the K3 cosection \(\mathrm{Ob}\twoheadrightarrow\cO\).  The orientation
line is fixed by the reduced K3 orientation and the product
Calabi--Yau form \(\sigma_S\wedge dz_E\).

\subsection*{Finite torsion descent}

The finite subgroup \(G=E[N]\cong(\bZ/N)^2\) acts freely on \(E\) and
trivially on \(S\).  The quotient
\[
  X/G=S\times E_N,\qquad E_N=E/E[N],
\]
is an isogenous copy of \(S\times E\).  A \(G\)-linearised coherent sheaf
on \(X\) descends to \(X/G\), and pullback from \(X/G\) produces a
\(G\)-linearised sheaf on \(X\).  This is finite etale descent, with
the usual character decomposition in the representation ring
\[
  R_G=\bZ[G^\vee].
\]

This descent is a finite equivariant descent datum.  In
cohomological Chow theory over \(\bQ\), \(H^\bullet(BG,\bQ)=\bQ\), so
finite torsion supplies no continuous localisation parameter.  In
equivariant \(K\)-theory, concentration after localising \(R_G\) gives
character idempotents and denominators \(\lambda_{-1}(N^\vee)\).  This
is \(K\)-theoretic fixed-point concentration.  The
Graber--Pandharipande virtual residue formula for the reduced DT class
requires a genuine \(\Gm\)-equivariant obstruction theory.

For the point-object charge this distinction is visible.  If
\[
  M_{[pt]}(X)=X,
\]
then every nonidentity element of \(E[N]\) translates the support of a
skyscraper sheaf.  Hence
\[
  M_{[pt]}(X)^G=\varnothing
\]
for nontrivial \(G\).  A point of \(X/G\) pulls back to the orbit-sum of
\(|G|=N^2\) points on \(X\); the quotient/descent lane begins at that
orbit charge.  The charge-one fixed residue is absent.

\subsection*{Finite K3 automorphism strata}

Finite symplectic K3 automorphisms give CHL-type strata only on special
K3 loci.  If \(g\in\mathrm{Aut}(S)\) has order \(N\) and preserves
\(\sigma_S\), then \(g\times t_a\) preserves
\(\sigma_S\wedge dz_E\) for every \(N\)-torsion translation \(t_a\) of
the elliptic factor.  Choosing \(a\) so that the product action is free
produces a smooth Calabi--Yau quotient.  If the quotient presentation is
kept stack-theoretic, the K3 fixed-point data enter through orbifold and
inertia sectors.  A global
\(\Gm\)-localisation parameter on \(S\times E\) is still absent, and
Nikulin's finite symplectic automorphism bounds control the allowable
orders.

\subsection*{K3-fibre reduced DT}

Let \(\pi:X\to E\) be the projection.  For a K3-primitive class
\(\beta\in H_2(S,\bZ)\), Pandharipande--Thomas prove the KKV
stable-pair formula on the K3 fibre:
\[
 Z^{\PT}_{S}(q)
 \;=\;
 \sum_{n}\chi(\Hilb^{n}(S))q^{n-1}
 \;=\;
 \prod_{k\geq 1}(1-q^{k})^{-24}.
\]
On \(X=S\times E\), the K3 symplectic form gives the cosection
\[
  \mathrm{Ob}\twoheadrightarrow\cO,
\]
and the reduced class has virtual dimension \(1\) before quotienting the
free elliptic translation.  Oberdieck--Pandharipande compute the
primitive reduced DT series as
\[
  Z^{X}_{\mathrm{OP/DT}}(q,t,p)
  =
  -(\PhiOP)^{-1}
  =
  -\Dfive^{-2}
  =
  -4096\,\Deltafive^{-2},
  \qquad
  \Dfive=64^{-1}\Deltafive,\qquad
  \PhiOP=\Dfive^2.
\]
Here \(\Deltafive\) is the primitive Gritsenko--Nikulin denominator of
weight \(5\), and \(\PhiOP\) is the Oberdieck--Pandharipande monic
weight-\(10\) scalar.  The variables \(q,t,p\) record elliptic base
degree, K3 curve degree, and point class in this convention.  The
primitive denominator and the scalar square carry different weights:
\[
  \kbkm(\Deltafive)=c_1(0)/2=10/2=5,\qquad
  \mathrm{wt}(\PhiOP)=10.
\]
Nonprimitive K3 classes require the imprimitive multiple-cover
correction.  The Oberdieck--Pixton programme controls broader
descendant and equivariant extensions; the \(N=1\) scalar displayed
above is the Oberdieck--Pandharipande reduced-DT theorem.

The \(d=3\) specialisation
\[
  \Phi_3^{(S,E)}\bigl(\Db\Coh(S\times E)\bigr)
\]
is an \(E_1\)-chiral algebra on the elliptic base.  The braided \(E_2\)
structure belongs to the Drinfeld-centre construction on
representations of this \(E_1\)-algebra.

\subsection*{Fibrewise Nakajima--Grojnowski theory}

Nakajima 1997 and Grojnowski 1996 construct a Heisenberg algebra
\(\cH(S)\) acting on
\[
  \bigoplus_{n\geq0}H^{*}(\Hilb^{n}(S),\bQ)
\]
for a smooth surface \(S\).  The creation operator \(p_{-k}(\alpha)\)
adds \(k\) points carrying \(\alpha\in H^{*}(S)\), and
\[
  [p_k(\alpha),p_{-\ell}(\beta)]
  =
  k\delta_{k,\ell}(\alpha,\beta)
\]
is the Heisenberg commutator.  The construction is a surface theorem.
For \(S\times E\), it applies fibrewise to
\[
  \Hilb^n_{S/E}(S\times E)=\Hilb^n(S)\times E,
\]
the relative Hilbert scheme of points contained in fibres of
\(\pi\).  The absolute Hilbert scheme of points on the threefold is
governed by DT/PT obstruction theory with three-dimensional
embedded-point deformations.

The fibre supplies the Mukai-enhanced Heisenberg algebra
\[
  H_{\mathrm{Muk}}(S),
  \qquad
  \operatorname{rank}\widetilde H^{*}(S,\bZ)=24.
\]
The elliptic direction supplies the spectral parameter \(z\in E\) by
evaluation modules
\[
  \mathrm{ev}_{z}\colon
  H_{\mathrm{Muk}}(S)_{\mathrm{aff}}\to H_{\mathrm{Muk}}(S).
\]
Thus the fibrewise Heisenberg currents and the elliptic spectral
parameter give the rank-\(24\) current block tested by the reduced DT
character.  The Borcherds bracket requires the denominator-root datum and
the Hall pairing; the surface Heisenberg action supplies the current
generators only.

\subsection*{Compatibility of the three mechanisms}

The three mechanisms carry different pieces of the \(K3\times E\) data.
Finite torsion descent supplies quotient and character data for free
\(E[N]\)-actions and for CHL-type finite automorphism strata.
Reduced DT supplies the Oberdieck--Pandharipande scalar
\(-\Dfive^{-2}\).  Fibrewise Nakajima--Grojnowski theory supplies the
\(24\) Mukai Heisenberg currents and their elliptic evaluation
parameter.

The full BKM-side quantum object is the Hall--Drinfeld double of the
compact oriented positive-half
\(\CoHA(K3\times E)\), once the negative half, Hopf pairing, Hall
coproduct, centre compatibility, completion, and Borcherds bracket
comparison are constructed.  The scalar reduced-DT identity proves a
positive-half character equality.  Double construction is the additional
Hall problem: build both halves, prove PBW, identify the radical quotient,
and compare the Serre/Borcherds relations with
\(\mathfrak g_{\Delta_5}\).

\subsection*{Point charge and fibre current checks}

At point-object charge, the moduli space is
\[
  M_{[pt]}=S\times E.
\]
For nontrivial \(G=E[N]\), the fixed locus \(M_{[pt]}^G\) is empty by
free translation on the elliptic factor.  The first quotient-descended
point class is the \(N^2\)-point orbit class.  Its ordinary topological
Euler diagnostic remains
\[
  \chi_{\mathrm{top}}(S\times E)=24\cdot0=0,
\]
consistent with
\[
  \kcat(S\times E)=\chi(\cO_{S\times E})
  =
  \chi(\cO_S)\chi(\cO_E)
  =
  2\cdot0
  =
  0.
\]
The K3 fibre value \(2=\chi(\cO_S)\) is a separate fibre witness.

The reduced K3 fibre contributes the \(24\)-dimensional Mukai
cohomology basis of the Heisenberg vacuum.  The Nakajima construction
realises the same basis as the creation operators
\[
  p_{-1}(e_a),\qquad a=1,\ldots,24,
\]
acting on \(H^{*}(\Hilb^1(S))=H^{*}(S)\).  Affinisation by evaluation
at \(z\in E\) gives currents \(p_{-1}(e_a;z)\).

The four construction-level invariants attached to the \(K3\times E\)
datum are
\[
  \{\kcat,\kch^{\mathrm{Heis}},\kbkm,\kfib\}(K3\times E)
  =
  \{0,3,5,24\}.
\]
They are read as
\[
  \kcat(K3\times E)=0,
  \qquad
  \kch^{\mathrm{Heis}}(K3\times E)=3,
\]
\[
  \kbkm(\Deltafive)=c_1(0)/2=10/2=5,
  \qquad
  \kfib(S)=\operatorname{rank}\widetilde H^{*}(S,\bZ)=24.
\]
The raw compact Hodge supertrace
\[
  \sum_q(-1)^qh^{0,q}(K3\times E)=0
\]
and the K3 structure-sheaf Euler characteristic \(\chi(\cO_S)=2\) are
separate witnesses.  The Heisenberg value \(3\), the Borcherds weight
\(5\), and the Mukai rank \(24\) come from distinct constructions.

\subsection*{Sources and remaining comparison}

Graber--Pandharipande 1999 supplies virtual localisation for
\(\Gm\)-equivariant obstruction theories.  Edidin--Graham supplies the
equivariant intersection framework, and finite \(E[N]\)-data here are
used as descent and character data.  The reduced DT input is the
Kiem--Li/Maulik--Pandharipande--Thomas cosection reduction with the
Oberdieck--Pandharipande orientation normalisation.
Oberdieck--Pandharipande prove the primitive \(K3\times E\)
reduced-DT scalar at \(N=1\); Oberdieck--Pixton belongs to the broader
descendant and equivariant extension.  Nakajima--Grojnowski supplies the
fibrewise Heisenberg algebra.

The compact Hall--Drinfeld double comparison remains the proof
obligation: construct the compact oriented \(\CoHA^{+}(K3\times E)\),
the negative half, Hopf pairing, coproduct, centre compatibility,
completion, and the Borcherds bracket/Serre comparison with
\(\mathfrak g_{\Delta_5}\).
